\begin{proposition}[Schouten像の三十八次元性の直接証明]
\label{prop:schouten-span-38-analytic-v31}
$X=\mathbb P(1,1,1,3)$ とする。標準座標とEulerベクトル場を用いて
\[
 H^0(X,\wedge^3T_X)\simeq
 R_6=\operatorname{Sym}^6\langle x_1,x_2,x_3\rangle
 \oplus w\operatorname{Sym}^3\langle x_1,x_2,x_3\rangle\oplus\mathbf Cw^2
\]
と同定すると、Schouten括弧の像は正確に
\[
 W=\operatorname{Sym}^6\langle x_1,x_2,x_3\rangle
   \oplus w\operatorname{Sym}^3\langle x_1,x_2,x_3\rangle
\]
である。特に $\dim W=28+10=38$ である。
\end{proposition}

\begin{proof}
$\Omega=dx_1\wedge dx_2\wedge dx_3\wedge dw$ とし、
重み零の三重ベクトル場 $U$ に対して
\[
 \Phi(U)=\Omega(E\wedge U)\in R_6
\]
と置く。この写像は $E\wedge\mathcal P^2$ 上で零であり、
射影的三重ベクトルの上記同定を与える。

重み零の二重ベクトルでは、$\partial_{x_i}\wedge\partial_{x_j}$ の係数は
$x$ の二次式であり、$\partial_{x_i}\wedge\partial_w$ の係数は
$x$ の四次式と $w$ に $x$ の一次式を掛けた式の和である。
二つの二重ベクトルの括弧を取ると、
$\partial_{x_1}\wedge\partial_{x_2}\wedge\partial_{x_3}$ の係数は $x$ の三次式、
$\partial_{x_i}\wedge\partial_{x_j}\wedge\partial_w$ の係数は
$x$ の五次式と $w$ に $x$ の二次式を掛けた式の和となる。
従って $\Phi$ の像に $w^2$ 項は生じず、Schouten像は $W$ に含まれる。

逆の包含を単項式ごとに示す。
$(i,j,k)$ を $(1,2,3)$ の巡回置換とし、
\[
 P=f\partial_{x_i}\wedge\partial_{x_j},\qquad
 Q=g\partial_{x_k}\wedge\partial_w,
 \qquad f=x^a,\quad g=w^\varepsilon x^b,
\]
\[
 |a|=2,\qquad |b|=d=4-3\varepsilon,\qquad \varepsilon\in\{0,1\}
\]
と置く。Jacobi式の混合項を直接計算すると、採用するSchoutenの符号規約による
共通の非零定数倍を除いて、
\begin{align*}
 \Phi([P,Q])
 &= -f(x_i\partial_{x_i}+x_j\partial_{x_j})g
    -gx_k\partial_{x_k}f\\
 &= (b_k-a_k-d)fg.
\end{align*}

まず $h=x^c$、$|c|=6$ を任意に取る。
$h=fg$ を $\deg f=2$、$\deg g=4$ と分解し、$c_k\le3$ となる $k$ を選ぶ。
そのような $k$ は必ず存在し、
\[
 b_k-a_k-4=c_k-2a_k-4\le-1
\]
だから、$h$ は括弧の像に入る。

次に $h=wx^c$、$|c|=3$ とする。
ある $c_k$ が零なら、その $k$ を選び、任意の
$h=fg$、$\deg_xf=2$、$g=w\cdot(\text{一次単項式})$ を使えば係数は $-1$ となる。
全ての $c_k$ が正なら $c=(1,1,1)$ である。
この場合は $f$ に現れる変数の一つを $x_k$ と選べば
$a_k=1$、$b_k=0$ なので係数は $-2$ となる。
従って $W$ の全単項式が括弧の像に入る。

ここで用いた $P,Q$ は発散零である必要はない。
発散零の代表へ変更した差はEulerベクトル場との外積であり、
その部分はSchoutenイデアルを成すため、射影的な括弧も $\Phi$ の値も変わらない。
\end{proof}
